\chapter{Avance del objetivo específico 2}

\section{Clasificador autosupervisado como módulo principal}

El clasificador es el núcleo del proyecto. El sistema usa un codificador DINOv2 ViT-B/14 público y estudia su adaptación mediante SimSiam con imágenes dermatológicas. Los brazos A* y B* corresponden a condiciones de adaptación distintas; B* incorpora el emparejamiento por tono explorado en el protocolo. La regresión logística posterior utiliza etiquetas de HAM10000 para evaluar la representación. Por tanto, el flujo completo contiene una fase SSL sin etiquetas y una etapa supervisada de evaluación de la tarea posterior.

\begin{figure}[htbp]
\centering
\begin{tikzpicture}[
    box/.style={draw, rounded corners, align=center, text width=2.4cm, minimum height=1.15cm, inner sep=4pt},
    main/.style={box, fill=blue!8},
    aux/.style={box, dashed, fill=orange!8},
    flow/.style={-{Stealth[length=2mm]}, thick},
    node distance=0.3cm
]
\node[main] (pool) {Imágenes sin etiqueta\\conjunto SSL versionado};
\node[main, right=of pool] (views) {Dos vistas\\aumentadas};
\node[main, right=of views] (encoder) {Codificador DINOv2 ViT-B/14\\+ adaptación SimSiam};
\node[main, right=of encoder] (probe) {Representación\\+ regresión logística};
\node[main, right=of probe] (classes) {Predicción\\de 7 clases};
\draw[flow] (pool) -- (views);
\draw[flow] (views) -- (encoder);
\draw[flow] (encoder) -- (probe);
\draw[flow] (probe) -- (classes);

\node[aux, below=0.9cm of encoder, text width=5.6cm] (segment) {Módulo complementario de segmentación\\MoCo v2/ResNet-50 $\rightarrow$ CCAM $\rightarrow$ pseudo-máscaras $\rightarrow$ U-Net};
\node[aux, right=0.8cm of segment, text width=4.1cm] (gate) {Compuerta S independiente\\Integración con clasificador: compuerta I, requiere ablación};
\draw[flow, dashed] (pool.south) |- (segment.west);
\draw[flow, dashed] (segment) -- (gate);
\end{tikzpicture}
\caption{Arquitectura modular del proyecto: clasificación SSL como núcleo y segmentación como componente independiente.}
\label{fig:arquitectura-modular}
\end{figure}
\noindent\textit{Nota.} La flecha entre segmentación y la compuerta I indica una evaluación posterior, no una integración ya validada. El diseño de preentrenamiento se apoya en DINOv2, SimSiam, MoCo v2 y CCAM \parencite{oquab2024dinov2,chen2021simsiam,chen2020mocov2,xie2022ccam}.

\section{Resultados disponibles de clasificación}

La tabla \ref{tab:v7-clasificacion} resume las medias de tres semillas de la evaluación v7. Los brazos adaptados presentan métricas puntuales superiores al DINOv2 público en HAM10000; los contrastes registrados favorecen de forma robusta a B* frente al público para el criterio P3-IC. En cambio, los contrastes B*--A* no sustentan una ventaja general del término de tono. La no inferioridad P1-NI de A* no es robusta porque falla en la semilla 42.

\begin{table}[htbp]
\centering
\caption{Resumen de clasificación v7 en HAM10000: medias de tres semillas.}
\label{tab:v7-clasificacion}
\begin{tabular}{lrrrr}
\toprule
Brazo & F1 macro & Exactitud balanceada & AUC macro & ECE \\
\midrule
DINOv2 público & 0,6158 & 0,5924 & 0,9257 & 0,1307 \\
A* sin emparejamiento por tono & 0,6719 & 0,6308 & 0,9640 & 0,0192 \\
B* con emparejamiento por tono & 0,6656 & 0,6284 & 0,9632 & 0,0236 \\
\bottomrule
\end{tabular}
\end{table}
\noindent\textit{Nota.} Elaboración propia a partir de los resultados v7 para las semillas 42, 43 y 44. Los intervalos de confianza publicados se remuestrearon por imagen; al existir imágenes repetidas por lesión, se requiere una reestimación agrupada por \texttt{lesion\_id} antes de tratarlos como inferencia confirmatoria.

La lectura secundaria v7.1 en validación reporta F1 macro de 0,7229 para A* y 0,7141 para B*, frente a 0,6320 para el control público con el protocolo de A*. A* y B* no se distinguen de forma consistente entre sí con el contraste pareado. La lectura v7.1 se seleccionó en la partición interna de selección y se evaluó posteriormente en validación; no sustituye el resultado primario v7 ni debe describirse como la primera lectura de validación de todo el proyecto.

\section{Avance y límites del módulo complementario de segmentación}

La línea de segmentación sin etiquetas sigue un flujo separado de MoCo v2/ResNet-50, CCAM, pseudo-máscaras y un U-Net estudiante. El candidato S4-A2 obtuvo Dice 0,7221 en ISIC2018 y 0,7949 en HAM10000; superó por márgenes pequeños algunos umbrales registrados frente al estudiante \texttt{student\_r3}. Sin embargo, las pruebas de integración C1 no muestran una mejora general del clasificador al usar máscaras. Este resultado respalda mantener la segmentación como módulo complementario, no afirmar que mejore el clasificador.

\noindent\textit{[Pendiente: insertar una figura representativa del segmentador y sus ejemplos de máscaras, con selección de casos y fuente verificadas.]}
